% Generated draft: requires mathematical and source review.
\documentclass[11pt]{article}
\usepackage[margin=25mm]{geometry}
\usepackage{fontspec}
\setmainfont{Latin Modern Roman}
\usepackage{amsmath,amssymb,mathtools,unicode-math}
\setmathfont{Latin Modern Math}
\usepackage{enumitem,needspace,xurl,hyperref}
\setlength{\parindent}{0pt}
\setlength{\parskip}{5pt}
\setlength{\emergencystretch}{3em}
\newcommand{\N}{\mathbb{N}}
\newcommand{\Z}{\mathbb{Z}}
\newcommand{\Q}{\mathbb{Q}}
\newcommand{\R}{\mathbb{R}}
\newcommand{\C}{\mathbb{C}}
\newcommand{\F}{\mathbb{F}}
\newcommand{\A}{\mathbb{A}}
\newcommand{\PP}{\mathbb{P}}
\newcommand{\E}{\mathbb{E}}
\newcommand{\eps}{\varepsilon}
\newcommand{\dd}{\,\mathrm{d}}
\begin{document}
\section*{UnsolvedMath ID: 30000827}
\textbf{Source problem code:} catalogue entry
\subsection{Definitions and mathematical statement}
For a path-connected space \(X\), use the unreduced convention: \(\operatorname{TC}(X)\) is the least number of open sets covering \(X\times X\) on each of which the endpoint map
\[
 \operatorname{ev}:X^{[0,1]}\to X\times X,\qquad
 \gamma\mapsto(\gamma(0),\gamma(1)),
\]
admits a continuous section. The reduced convention subtracts one. The record does not specify which normalization it uses.

Let \(n\geq3\) and \(\ell=(\ell_1,\ldots,\ell_n)\in\R_{>0}^n\). Under the usual genericity convention, \(\sum_i\varepsilon_i\ell_i\ne0\) for every \(\varepsilon_i\in\{1,-1\}\). The labelled planar polygon space modulo translations and rotations can be written
\[
 M_\ell=\left\{(u_1,\ldots,u_n)\in(S^1)^n:
                  \sum_{i=1}^n\ell_i u_i=0\right\}/\operatorname{SO}(2),
\]
where rotations act diagonally. Reflections are not included in this quotient, nor are arbitrary permutations of edges; self-intersection is not excluded by the supplied statement. The first request is to determine \(\operatorname{TC}(M_\ell)\) for arbitrary generic length data in the intended nonempty, connected domain. Nonemptiness requires the polygon inequalities, and genericity alone does not ensure that the space is nonempty or path-connected. With the endpoint-map definition on all of \(X\times X\), a space with multiple path components has no finite topological complexity because some pairs cannot be joined by a path; restricting to components or also quotienting by reflections would be different problems.

The second question asks whether, for closed connected three-manifolds \(M,N\),
\[
 \pi_1(M)\cong\pi_1(N)\quad\Longrightarrow\quad
 \operatorname{TC}(M)=\operatorname{TC}(N).
\]
No orientability condition is explicitly supplied. In particular, for \(p\geq2\) and \(\gcd(p,q)=1\), write
\[
 L(p,q)=S^3/\bigl((z_1,z_2)\sim(\zeta z_1,\zeta^qz_2)\bigr),
 \qquad \zeta=e^{2\pi i/p}.
\]
Since these lens spaces have fundamental group \(\Z/p\Z\), the explicit special question is
\[
 \operatorname{TC}(L(p,q))\stackrel{?}{=}\operatorname{TC}(L(p,1))
 \quad\text{for all such }p,q.
\]
The lens-space comparison is a necessary instance of the broader group-dependence assertion, not an asserted equivalent reformulation.
\subsection{Short English statement}
The record asks for the motion-planning complexity of generic planar polygon spaces and whether that complexity for a closed three-manifold is determined by its fundamental group. Lens spaces with the same cyclic fundamental group provide the stated test case. Normalization and connectedness conventions matter for the polygon question.
\subsection{Sources}
\begin{itemize}
\item Local UnsolvedMath registry, permanent ID \(30000827\), exact source code \texttt{catalogue entry}, proposed year 2007; no individual proposer is supplied.
\item \emph{Mini-Workshop: Topology of closed one-forms and Cohomology Jumping Loci (2008)}, Oberwolfach Reports, DOI 10.4171/owr/2007/40, pp. 2321--2360. \url{https://doi.org/10.4171/owr/2007/40}.
\item Supplied reference labelled ``planar polygon spaces'': \url{https://arxiv.org/abs/1504.00848}.
\item The supplied \texttt{corrected\_verified} note reports checking mathematical content against the official report. The assessment checked 2026-08-21 reports computations for important polygon-space families and sharp alternatives in broad classes, while arbitrary length data and the general three-manifold group-dependence question remain unsettled. This and the partially solved label are dated supplied triage requiring independent verification.
\end{itemize}
\end{document}
